\section{Introdução}
\label{sec:introducao}

O desenvolvimento de veículos autônomos tem despertado crescente interesse da indústria automobilística e da comunidade científica. No setor automotivo, esse interesse está associado à expansão do mercado de sistemas avançados de assistência ao condutor (\textit{Advanced Driver Assistance Systems} -- ADAS) e às perspectivas de desenvolvimento e comercialização de tecnologias com níveis mais elevados de automação. Segundo estimativas da McKinsey, o mercado global de software e componentes eletrônicos relacionados a essas tecnologias poderá alcançar aproximadamente US\$ 160 bilhões até 2035, com crescimento anual projetado de cerca de 16\% \citep{McKinsey2026}.

No âmbito científico, as pesquisas são motivadas pelo potencial dessas tecnologias para aumentar a segurança e a eficiência dos sistemas de transporte, bem como limitações nas metodologias de teste desses sistemas, particularmente na correspondência entre os resultados obtidos em simulação e o comportamento observado em ambientes reais \citep{Liao2025}.

No Brasil, diferentes grupos de pesquisa têm desenvolvido plataformas experimentais voltadas ao estudo de tecnologias de condução autônoma. Entre essas iniciativas, destacam-se os projetos CADU, da Universidade Federal de Minas Gerais (UFMG) \cite{Sabbagh2010}, CaRINA, da Universidade de São Paulo (USP) \cite{Fernandes2014}, e IARA, da Universidade Federal do Espírito Santo (UFES) \cite{Badue2021}. Essas iniciativas evidenciam o interesse da comunidade científica nacional pelo desenvolvimento e pela validação experimental de tecnologias aplicadas a veículos autônomos. 

Nesse contexto, o Laboratório de Mobilidade Terrestre (LMT), da Universidade Federal de Lavras (UFLA), iniciou em 2019 o desenvolvimento da plataforma VIDA (Veículo Inteligente de Desenvolvimento Aplicado). Trata-se de um veículo elétrico experimental concebido para apoiar pesquisas em automação e navegação de veículos terrestres \citep{Silva2024CIUFLA}.

Entre os desafios da condução autônoma, destaca-se o seguimento de trajetórias, que consiste em determinar as ações de controle necessárias para conduzir um veículo ao longo de um percurso de referência. Para isso, é necessária a integração entre os sistemas de localização e controle do veículo.

Diante desse cenário, este trabalho tem como objetivo implementar e validar experimentalmente um sistema de seguimento autônomo de trajetórias na plataforma VIDA, utilizando localização por fusão sensorial de encoders, sensor de esterçamento, giroscópio e receptor GNSS com correções RTK e controle lateral baseado no método Stanley.

A implementação envolve também o controle de baixo nível e uma arquitetura computacional distribuída baseada no ROS 2 (\textit{Robot Operating System 2}), composta por um microcontrolador ESP32, um computador embarcado Odroid e um notebook responsável por tarefas de alto nível. A validação é realizada mediante ensaios com o veículo real, considerando o acompanhamento da rota, os erros lateral e angular e o desempenho dos sistemas de controle.

Busca-se, assim, contribuir para a consolidação do VIDA como plataforma experimental para pesquisas futuras em veículos autônomos. Além da contribuição tecnológica e científica, o desenvolvimento deste trabalho também representa uma etapa importante na formação acadêmica em Engenharia de Controle e Automação, ao proporcionar a aplicação prática e integrada dos conhecimentos adquiridos ao longo da graduação.

Este artigo está estruturado da seguinte maneira: a Seção~\ref{sec:fundamentacao} apresenta a fundamentação teórica, enquanto a Seção~\ref{sec:plataforma} descreve a plataforma VIDA e seus principais subsistemas. Na Seção~\ref{sec:arquitetura}, é detalhada a arquitetura computacional adotada, seguida pela descrição dos controladores de baixo nível na Seção~\ref{sec:baixo_nivel}. As Seções~\ref{sec:estimacao} e~\ref{sec:seguimento_trajetoria} abordam, respectivamente, a estimação de estados e a implementação do seguimento de trajetórias. Por fim, a Seção~\ref{sec:resultados_experimentais} apresenta os resultados experimentais, e a Seção~\ref{sec:conclusao} reúne as conclusões do trabalho. 